\section{A line-potential certificate for the true-trail kernel}
\label{sec:h3line}

This section gives a finite replacement for the stationary-square input
in Appendix C.  Distances are lattice graph distances.  An unoriented
bond is an unordered pair of nearest-neighbour vertices; a trail never
uses the same such bond twice.  Revisiting a site means an earlier
\emph{departure} from it.  In particular, the incoming bond at the start
is not a used bond.  All sums below have nonnegative terms; an infinite
sum means the supremum of its finite partial sums until finiteness is
proved.  The result concerns the plane true-trail kernel, with the centre
excluded and the target stopped at its first hit.

The constants and the orderings in the present section replace those in
the kernel-order definition of Section 7.  The local thresholds, the
response coefficients, the meaning of a trail, and the kernel definition
itself are unchanged.  No change of the local input (H1) is needed for
this replacement.

\subsection{The kernel and the finite state space}

Encode the four absolute directions by
\[
 0=\mathrm E=(1,0),\quad 1=\mathrm N=(0,1),\quad
 2=\mathrm W=(-1,0),\quad 3=\mathrm S=(0,-1).
\]
Direction addition is modulo four; the opposite of $a$ is $a+2$.
Set
\begin{equation}\label{hl:constants}
\begin{gathered}
 x_{\rm s}=\frac{1000}{961},\qquad x_{\rm t}=\frac{200}{187},\qquad
 h_B=1,\quad h_C=\frac65,\quad \rho_0=\frac{967725}{10^6}
             =\frac{38709}{40000},\\
\begin{array}{c|ccc}
 \tau&a_\tau&b_\tau&b'_\tau\\\hline
 \mathrm s&145217/10^6&352782/10^6&340008/10^6\\
 \mathrm t&150000/10^6&357549/10^6&344661/10^6
\end{array}
\end{gathered}
\end{equation}
These are the rational ceilings justified by the logarithmic-mean
lemma in Section 7; in particular $0<b'_\tau\le b_\tau$.
For distinct directions $e,e'$ write
\[
 H(e,e')=\begin{cases}1,&e'=e+2,\\6/5,&e'\ne e+2.\end{cases}
\]
Thus $H$ is the $h$-weight of the pair of neighbours in these directions.

For row $X=B,C$ fix $c=(0,0)$, $u=(-1,0)$, and respectively
$v=(1,0)$ or $v=(0,1)$.  Let $\mathcal T_n^X$ be the length-$n$
trails from $u$, avoiding $c$, ending at their first hit of $v$.
The incoming direction at $u$ is W; at other departures it is the
preceding step.  At departure $i$ let $N_i(y)$ count departures from
$y$ at indices strictly less than $i$.  A neighbour is dead at that
time if it is $c$ or $N_i(y)\ge2$.  A first departure is one with
$N_i(v_i)=0$.  For a first departure with incoming $a$, outgoing $d$,
and $\{f,g\}=\{0,1,2,3\}\setminus\{a+2,d\}$, put $\tau=\mathrm s$
if $d=a$, otherwise $\tau=\mathrm t$.  Give the pair $\{f,g\}$
coefficient $a_\tau$ and each of $\{f,d\},\{g,d\}$ coefficient
$b_\tau$ when both $f,g$ are live.  When only one of $f,g$ is live,
give its pair with $d$ coefficient $b'_\tau$.  When neither is live
give no pair.  There is no response at a revisit.  The notation for a
pair of directions also denotes its pair of neighbour vertices.
For a live pair $p$ put
\[
 s(p)=\max_{z\in p}\min\{|z-u|_1,|z-v|_1\},
\]
and let $\kappa_{Y,s}(i)$ sum its coefficients of type $Y$ and shift
$s$.  Put $\kappa_i=\sum_{Y,s}\kappa_{Y,s}(i)h_Y$.

The outgoing neighbour is live on any such trail.  A noninitial site
departed from twice has used two arrival and two departure bonds, hence
all four incident bonds; it cannot be entered again.  At $u$, two
departures have used its three bonds other than $\{u,c\}$, also
precluding a further arrival.  The target has no previous departure.
This argument also applies to every centre-avoiding prefix generated
below, before any future completion is required.

Fix $L=12$.  For $n\le12$ order first departures by increasing
\begin{equation}\label{hl:order}
 \frac{x_i-1}{x_i\kappa_i},\qquad x_i=x_{\tau_i},
\end{equation}
using $+\infty$ when $\kappa_i=0$ and increasing path index to break
ties.  For $n>12$ order straight first departures in path order,
then turning first departures in path order.  Write this fixed,
$\theta$-independent ordering as $\sigma_w$.  The kernel is
\begin{equation}\label{hl:kernel}
\begin{split}
 M_{XY,s}&=\sum_{n\ge1}\sum_{w\in\mathcal T_n^X}
   3^{-n}(3/5)^{\operatorname{rev}(w)}
   \sum_k\kappa_{Y,s}(\sigma_w(k))
                      \prod_{l\le k}x_{\sigma_w(l)},\\
 M_\theta[X,Y]&=\sum_{s\ge0}M_{XY,s}\theta^{-s}.
\end{split}
\end{equation}
Here $\operatorname{rev}(w)$ counts revisit departures.  In particular
$3^{-n}(3/5)^{\operatorname{rev}(w)}$ gives coefficient $1/3$ at
each first departure and $1/5$ at each revisit.

The memory state space is obtained as follows.  Let $\Omega_{10}$ be
the direction words $\omega=(d_1,\ldots,d_{10})$ whose path
$z_0=0,z_j=\sum_{i\le j}d_i$ has distinct unoriented bonds.
At its endpoint, $d$ is admitted precisely when $d\ne d_{10}+2$
and $\{z_{10},z_{10}+d\}$ is not one of these bonds.  Its successor
is $(d_2,\ldots,d_{10},d)$ and its kind is
\[
 k(\omega,d)=\begin{cases}
 \mathrm r,&z_{10}\in\{z_0,\ldots,z_9\},\\
 \mathrm s,&z_{10}\notin\{z_0,\ldots,z_9\},\ d=d_{10},\\
 \mathrm t,&\text{otherwise}.
 \end{cases}
\]
The successor is in $\Omega_{10}$ because its bonds are a subset
of the distinct bonds of the admitted extended word, translated.

The quotient data are a map $\pi:\Omega_{10}\to Q$, with
$Q=\{0,\ldots,3740\}$, a representative $\omega_q$ for each
class, and four labelled successor/kind entries per class.
The following are part of the finite assertion, not assumptions
about a minimization algorithm.  Enumerate all words by trying all
four directions recursively, rejecting exactly a previously used
unoriented bond.  At depths $1,\ldots,10$ their counts must be
\begin{equation}\label{hl:windowcounts}
 4,12,36,108,316,916,2628,7500,21268,60092.
\end{equation}
Require that the map lists exactly these depth-ten words in
lexicographic order, with codes
$e(\omega)=\sum_{i=1}^{10}d_i4^{10-i}$, and valid class numbers.
Require $\pi(\omega_q)=q$.  For \emph{every} raw word and direction
require that admission and kind agree with the stored class row,
and that an admitted successor has the stored image under $\pi$.
For every admitted non-revisit direction also require that the raw
word and its representative have the same final direction.
Inadmissibility is recorded by successor and kind both $-1$;
kind codes $0,1,2$ mean s,t,r.  Require exactly $2516,5032,1872$
admitted class directions of these three kinds.  Write $q d$ for
the admitted successor and $k(q,d)$ for its kind.

These tests imply all the needed quotient properties.  Representatives
give surjectivity.  The coefficient of a non-revisit edge can use
$a_q$, the final direction of $\omega_q$, independently of the
raw representative.  A class with several possible final directions
has only revisit edges, where no such direction is used.  For any
function of phase, class and position, its pullback by $\pi$ has
exactly the same backward row as the raw operator: every direction
and next phase has the same coefficient, position and successor
image.  Directions remain distinct summation indices even if
their successor classes coincide.  There is no class-size factor.
The same identity for finite stopped sums follows by induction.

\subsection{The free and stationary operators}

The four phases are A,B,C,D, in this order.  The diagonal weights
$d_p(k)$ of an admitted direction are
\begin{equation}\label{hl:diagonals}
\begin{array}{c|ccc}
 p\backslash k&\mathrm s&\mathrm t&\mathrm r\\\hline
 \mathrm A&x_{\rm s}/3&x_{\rm t}/3&1/5\\
 \mathrm B&x_{\rm s}/3&1/3&1/5\\
 \mathrm C&x_{\rm s}/3&1/3&1/5\\
 \mathrm D&1/3&1/3&1/5
\end{array}.
\end{equation}
In addition there can be a change A$\to$C on a turning first
departure, or B$\to$D on a straight first departure.  Its coefficient
is $x_\tau K/3$.  No phase change is allowed at a revisit.  For
non-revisit incoming $a$, outgoing $d$, the free reward is
\begin{equation}\label{hl:freereward}
 K^\circ(a,d)=a_\tau H(f,g)+b_\tau\{H(f,d)+H(g,d)\}.
\end{equation}
In particular it is
$a_{\rm s}+2b_{\rm s}(6/5)$ on a straight passage and
$b_{\rm t}+(a_{\rm t}+b_{\rm t})(6/5)$ on a turn.

An individual edge $e$ specifies an absolute direction and a next
phase.  Its destination phase and class will be denoted $p_e,q_e$,
its displacement $d_e$, and its exact rational coefficient $t_e$.
The free operator, including its position displacement, is
\begin{equation}\label{hl:free}
 (\mathcal W H)(p,q,x)=\sum_{e\text{ free from }(p,q)}
                         t_e H(p_e,q_e,x+d_e).
\end{equation}
On position-independent functions this is the free operator
$\mathsf W$ in Appendix C.  All three response neighbours are live
in this definition.  The diagonal edge and change-of-phase edge
are separate; all pair contributions to one change-of-phase edge
are summed \emph{before} rounding its coefficient.

For a case $(X,d_0)$ with $d_0\in\{\mathrm N,\mathrm W,\mathrm S\}$,
put $a=u+d_0$.  Define $T=T_{X,d_0}$ on phases, $Q$, and
$\mathbb Z^2$ by the following permanent rules.
At $c$ and $v$ the row is zero.  At $u$ and $a$, forbid either
direction using $\{u,a\}$, kill every memory-detected revisit
departure, give every remaining admitted direction weight $1/5$
in the same phase, and give no response edge.  At other positions
use \eqref{hl:diagonals}; for the response declare just $c$ dead.
More explicitly, if $c$ is a neighbour in direction $\delta$, use
that $\delta$, otherwise use $\delta=4$.  For a direction not
entering $c$ its response reward is
\begin{equation}\label{hl:stationaryreward}
 K(a_q,d,\delta)=\begin{cases}
 b'_\tau H(g,d),&\delta=f,\\
 b'_\tau H(f,d),&\delta=g,\\
 a_\tau H(f,g)+b_\tau\{H(f,d)+H(g,d)\},&\delta\notin\{f,g\}.
 \end{cases}
\end{equation}
The response coefficient is again $x_\tau K/3$ on the indicated
phase change.  Directions entering $c$ have zero contribution and
may be deleted.  Keeping them with a zero continuation gives the
same patch rows.  The reward function on the state space is
\begin{equation}\label{hl:g}
  g(p,q,x)=\mathbf1_{\{x=v\}}g_p,\qquad(g_{\rm A},g_{\rm B},g_{\rm C},g_{\rm D})=(0,0,1,1).
\end{equation}
In particular $c$ is a zero-reward terminal.  These rules are never
released at a later suffix time.

For every nonnegative $H$,
\begin{equation}\label{hl:entrywise}
 TH\le\mathcal WH
\end{equation}
pointwise, matching each retained direction and next phase.
Indeed the usual diagonal coefficients agree.  Replacing one
dead non-outgoing neighbour by a live one increases a sole surviving
pair coefficient from $b'_\tau$ to $b_\tau$ and adds only
nonnegative pair terms.  At a known site $1/5$ is at most every
free diagonal coefficient, all response edges disappear, and more
directions can be killed.  Terminal rows and edges into $c$ only
decrease the operator.  This comparison retains positions; it is
valid on position-dependent functions as well as on constants.

\subsection{The exact finite assertion H3-line}

Fix
\begin{equation}\label{hl:scales}
 m=10,\quad L=12,\quad R=70,\quad B=20,\quad
 V=2^{20},\quad D=2^{48},\quad r=12/13.
\end{equation}
Here $B$ denotes the patch radius, not a pair or phase type.  For
any rational $t=A/B_0\ge0$ with integer $B_0>0$ define
\begin{equation}\label{hl:rounding}
 P(t)=\left\lceil Dt\right\rceil
 =\left\lfloor\frac{DA+B_0-1}{B_0}\right\rfloor.
\end{equation}
The defining integer enclosure is $DA\le B_0P(t)$ and, for
$A>0$, $B_0(P(t)-1)<DA$; for $A=0$ use $P(t)=0$.
Thus $t\le P(t)/D$.  Each coefficient is regenerated from
\eqref{hl:constants}--\eqref{hl:stationaryreward}.  For reference,
the four diagonal ceiling numerators, in order
$x_{\rm s}/3,x_{\rm t}/3,1/3,1/5$, are
\[
 97632666219444,\quad100347585280092,\quad
 93824992236886,\quad56294995342132.
\]

\begin{definition}[H3-line]\label{hl:finite}
The quotient checks above and all of the following finite
equalities and inequalities are required.

\smallskip\noindent
\textup{(i) Tilted vectors.}
Supply positive integers $Z^\sigma_{p,q}$ for
$\sigma\in\{+e_1,-e_1,+e_2,-e_2\}$, all phases $p$ and $q\in Q$,
and a positive integer $\eta$.  Require, for every such triple,
\begin{equation}\label{hl:tiltcheck}
 \sum_{e\text{ free from }(p,q)}P(t_e r^{\sigma\cdot d_e})
                         Z^\sigma_{p_e,q_e}+\eta
             \le D Z^\sigma_{p,q},
 \qquad Z^\sigma_{\mathrm C,q},Z^\sigma_{\mathrm D,q}\ge V.
\end{equation}
The tilt coefficients are exactly rational: $r^{-1}=13/12$,
$r^0=1$, $r^1=12/13$; round the \emph{product} $t_e r^{\sigma\cdot d_e}$
upward using \eqref{hl:rounding}.  Do not tilt a previously rounded
coefficient.  Require also the finite tail comparison
\begin{equation}\label{hl:tailbound}
 Z_{\max}12^{71}\le V13^{71},\qquad
 Z_{\max}=\max_{\sigma,p,q}Z^\sigma_{p,q}.
\end{equation}

\smallskip\noindent
\textup{(ii) Lines, caps, and both interfaces.}
For axis $j=1,2$, supply nonnegative integers
$E^j_{p,q,k}$ for $-70\le k\le70$.  Define the exact rational
cap numerator and its two finite roundings by
\begin{equation}\label{hl:cap}
 C^j_{p,q}(k)=\min\{Z^{+e_j}_{p,q}r^k,Z^{-e_j}_{p,q}r^{-k}\},
 \quad A^j_{p,q,k}=\lfloor C^j_{p,q}(k)\rfloor\ (|k|\le70),
 \quad B^j_{p,q,\pm}=\lceil C^j_{p,q}(\pm71)\rceil.
\end{equation}
These are derived numbers.  If stored, compare every one with
the displayed rational formula.  Require
\begin{equation}\label{hl:linecaps}
 \begin{gathered}
 0\le E^j_{p,q,k}\le A^j_{p,q,k}\qquad(|k|\le70),\\
 E^j_{\mathrm C,q,0}=E^j_{\mathrm D,q,0}=V,\qquad
 E^j_{\mathrm C,q,k},E^j_{\mathrm D,q,k}\le V\qquad(0<|k|\le70).
 \end{gathered}
\end{equation}
For the finite line rows set
\[
 \widehat E^j_{p,q,k}=
 \begin{cases}
 E^j_{p,q,k},&|k|\le70,\\
 B^j_{p,q,-},&k=-71,\\ B^j_{p,q,+},&k=71.
 \end{cases}
\]
For every $p,q,k$ with $|k|\le70$, except $p\in\{\mathrm C,\mathrm D\}$
and $k=0$, require
\begin{equation}\label{hl:linerow}
 \sum_{e\text{ free from }(p,q)}P(t_e)
             \widehat E^j_{p_e,q_e,k+(d_e)_j}
                    \le D E^j_{p,q,k}.
\end{equation}
Thus all A,B rows at zero and both endpoint rows $k=\pm70$ are
checked.  The cap inequalities at $k=\pm70$ control rows immediately
\emph{outside} the stored interval; the cap ceilings at $\pm71$
control rows leaving it.  Requiring the caps at every interior $k$
is a harmless strengthening matching the supplied data.

The mathematical line potential is defined without further rounding:
\begin{equation}\label{hl:ell}
 \ell^j_{p,q}(k)=\begin{cases}
 E^j_{p,q,k}/V,&|k|\le70,\\
 C^j_{p,q}(k)/V,&|k|>70.
 \end{cases}
\end{equation}
For a target $v$ put
\begin{equation}\label{hl:Finfty}
 F^v_\infty(p,q,x)=
     \min\{\ell^1_{p,q}(x_1-v_1),\ell^2_{p,q}(x_2-v_2)\}.
\end{equation}
The same $Z,E$ are shared by all five cases.

\smallskip\noindent
\textup{(iii) Five patches.}
For $(X,d_0)=(B,\mathrm N),(B,\mathrm W),(C,\mathrm N),
(C,\mathrm W),(C,\mathrm S)$, supply integers
$U^{X,d_0}_{p,q,x}$ on $P=[-20,20]^2\cap\mathbb Z^2$.
Require, for every index,
\begin{equation}\label{hl:patchcaps}
 0\le U_{p,q,x}\le V F^v_\infty(p,q,x),\qquad
 U_{p,q,c}=0,\qquad U_{p,q,v}=Vg_p.
\end{equation}
There is no outgoing row to check at $c$ or $v$.  Define for a
successor $y$ of a site in $P$
\[
 \widehat U_{p,q,y}=\begin{cases}
 U_{p,q,y},&y\in P,\\
 V F^v_\infty(p,q,y),&y\notin P.
 \end{cases}
\]
All coordinates $y_j-v_j$ in this formula have absolute value at
most $22<70$.  Consequently its outside value is the integer
$\min\{E^1_{p,q,y_1-v_1},E^2_{p,q,y_2-v_2}\}$, with no rounding.
Require every stationary nonterminal row
\begin{equation}\label{hl:patchrow}
 \sum_{e\text{ stationary from }(p,q,x)}P(t_e)
                  \widehat U_{p_e,q_e,x+d_e}
                    \le D U_{p,q,x},\qquad x\in P\setminus\{c,v\}.
\end{equation}
In particular this includes every boundary site of $P$ and every
edge leaving it, and uses the permanent known-site rules above.

\smallskip\noindent
\textup{(iv) Exact prefixes.}
Define a full-lattice function for each case by
\begin{equation}\label{hl:Fpatch}
 F(p,q,x)=\begin{cases}U_{p,q,x}/V,&x\in P,\\F^v_\infty(p,q,x),&x\notin P.\end{cases}
\end{equation}
Let $S_{X,d_0}$ be the exact sum of the $h$-weighted expressions
in \eqref{hl:kernel} over first-hit trails with $1\le n\le12$
and first direction $d_0$, using \eqref{hl:order}.
Let $\mathcal P_{X,d_0}$ consist of all length-twelve true prefixes
from $u$, avoiding both $c$ and $v$, with that first direction.
For such a prefix, propagate a phase vector $(p_{\rm A},p_{\rm B},p_{\rm C},p_{\rm D})$
starting at $(1,1,0,0)$, using its actual revisits, actual dead
neighbours and actual reward $\kappa$, as follows:
\begin{equation}\label{hl:phaseupdate}
\begin{array}{ll}
 \mathrm r:&p\longmapsto p/5,\\[2pt]
 \mathrm s:&p\longmapsto
 (x_{\rm s}p_{\rm A}/3,\ x_{\rm s}p_{\rm B}/3,\ x_{\rm s}p_{\rm C}/3,
             (p_{\rm D}+x_{\rm s}\kappa p_{\rm B})/3),\\[2pt]
 \mathrm t:&p\longmapsto
 (x_{\rm t}p_{\rm A}/3,\ p_{\rm B}/3,
             (p_{\rm C}+x_{\rm t}\kappa p_{\rm A})/3,\ p_{\rm D}/3).
\end{array}
\end{equation}
All right sides use the old components.  The phase weights in a
prefix use the straight-first construction even if its length is
at most twelve.  For a \emph{short hit}, use the sorted kernel
expression, not its phase-vector terminal reward.  If $q(P)$ is
the image of the last ten directions and $x(P)$ the endpoint, set
\begin{equation}\label{hl:J}
 J_{X,d_0}=S_{X,d_0}+\sum_{P\in\mathcal P_{X,d_0}}\sum_p
                    p_p(P)F(p,q(P),x(P)).
\end{equation}
Require the five exact inequalities
\begin{equation}\label{hl:prefixceil}
\begin{array}{c|rrrrr}
 (X,d_0)&B/\mathrm N&B/\mathrm W&C/\mathrm N&C/\mathrm W&C/\mathrm S\\\hline
 10^9 J_{X,d_0}\le&343236463&281251920&522717990&372660908&264538742.
\end{array}
\end{equation}
There are no floating-point comparisons in this assertion.
\end{definition}

Here is a precise integer interpretation of part (iv), including
an enumeration contract.  Start a recursive node at
$j=0,x=u$, incoming W, no used bonds, and departure counts zero.
Try directions in E,N,W,S order, retaining just $d_0$ at $j=0$.
Reject a move to $c$, a used unoriented bond, and the initial
reverse direction.  For each admitted move, compute the revisit
flag and all dead-neighbour flags from the predeparture counts,
then increment the count at $x$, add its bond, move, and recurse
with a copy of the data.  Keep each first departure's
$(i,x_i,\kappa_i)$ recorded at its departure time.  Test $x=v$
\emph{before} $j=12$: a hit, including a hit at twelve, contributes
its sorted short value and has no children.  A node at $j=12$
that is not a hit contributes its phase dot product in
\eqref{hl:J} and has no children.  All other nodes continue.
The induction proving completeness is elementary: every retained
word has the stated used bonds and departure counts; conversely
each permitted word has a unique parent and passes precisely
these tests.  This proves both exhaustion and absence of duplicates.

For this $h$, a common denominator, in fact the least common
multiple of the relevant reduced edge denominators, is
\begin{equation}\label{hl:Delta}
 \Delta=13478025000.
\end{equation}
This is checked by regenerating $1/3,1/5,x_{\rm s}/3,x_{\rm t}/3$
and $x_\tau\kappa/3$ for all nonreversing incoming/outgoing
directions and all four choices of dead flags on $f,g$, using
\eqref{hl:constants}.  Require divisibility of $\Delta$ by
every reduced denominator; the least-common-multiple equality
can also be checked and is used by the replay.  In units of
$\Delta^{-j}$ the phase vector is integral.  A revisit multiplies
each numerator by $\Delta/5$.  For a straight departure, if
$A,B,C,D_0$ denote the four old numerators, put
\[
 (A,B,C,D_0)\longmapsto
 (n_{\rm s}A,n_{\rm s}B,n_{\rm s}C,n_3D_0+n_\kappa B),
 \quad
 n_{\rm s}=\Delta x_{\rm s}/3,\quad n_3=\Delta/3,\quad
 n_\kappa=\Delta x_{\rm s}\kappa/3;
\]
for a turn replace this by
$(n_{\rm t}A,n_3B,n_3C+n_\kappa A,n_3D_0)$, with the analogous
$n_{\rm t},n_\kappa$.  All these integers are regenerated exactly.

Every twelve-step endpoint has $|x-u|_1\le12$, hence lies in
$[-13,11]\times[-12,12]\subset P$.  It also has
$x_1+1+x_2\equiv0\pmod2$.  If $A_{X,d_0}$ denotes the sum of
the four resulting phase numerators times the corresponding
stored $U$ values, and $S_{X,d_0}=s/d$ with $d>0$, then
\eqref{hl:prefixceil} with bound $K$ means exactly
\begin{equation}\label{hl:clearedprefix}
 10^9\{s\Delta^{12}V+A_{X,d_0}d\}
                       \le Kd\Delta^{12}V.
\end{equation}
The endpoint is looked up in the full patch.  Any compact lookup
must agree entry by entry with that full lookup for every used
phase, class, and coordinate.  Parity is only an index restriction,
and introduces no factor of two.

For completeness, the short sorting also has an exact integer
interpretation: for positive rewards compare
$(x_i-1)x_j\kappa_j$ and $(x_j-1)x_i\kappa_i$ by clearing
positive rational denominators, and use path index on equality.
With preceding product $P_0$, the difference between the
contributions in orders $i,j$ and $j,i$ is
\begin{equation}\label{hl:exchange}
 P_0\{x_j\kappa_j(x_i-1)-x_i\kappa_i(x_j-1)\}.
\end{equation}
The product after both entries agrees.  Thus increasing
\eqref{hl:order} minimizes the expression by adjacent exchanges;
an equal-key exchange changes no value.  A zero-reward entry is
placed last because moving it past a positive-reward entry decreases
the value.  None of this uses a dyadic approximation to $6/5$.
The optimality observation explains the ordering, while the proof
below only needs that this is the exact ordering defining $M$.

\subsection{From the finite line tests to an infinite potential}

\begin{lemma}[Exact tails and interfaces]\label{hl:taillemma}
For each axis $j$, the line potentials of \eqref{hl:ell} are
nonnegative and finite-valued.  They satisfy
\begin{equation}\label{hl:lineexact}
 \sum_{e\text{ free from }(p,q)}t_e
             \ell^j_{p_e,q_e}(k+(d_e)_j)\le\ell^j_{p,q}(k)
\end{equation}
for all $k\in\mathbb Z$ when $p$ is A or B, and for every
$k\ne0$ when $p$ is C or D.  In the latter two phases the values
are at most one for every $k$ and equal one at $k=0$.
\end{lemma}

\begin{proof}
First use the exact exponentials on the whole integer line,
\[
 E^+_{p,q}(k)=Z^{+e_j}_{p,q}r^k/V,\qquad
 E^-_{p,q}(k)=Z^{-e_j}_{p,q}r^{-k}/V.
\]
Divide \eqref{hl:tiltcheck} by $DV$ and use
$t_e r^{\pm(d_e)_j}\le P(t_e r^{\pm(d_e)_j})/D$.  Multiplying
by $r^{\pm k}$ gives, for either sign and every integer $k$,
\[
 \sum_e t_e E^\pm_{p_e,q_e}(k+(d_e)_j)
 \le E^\pm_{p,q}(k)-\frac{\eta r^{\pm k}}{DV}
 \le E^\pm_{p,q}(k).
\]
All signs are rational powers, so no limiting or rounding error
is involved.  The positive residual is more than is needed here.

For a nonnegative linear row operator $S$ and functions $f,g$,
monotonicity gives
\begin{equation}\label{hl:min}
 S(\min(f,g))\le Sf,\quad S(\min(f,g))\le Sg,
 \quad\text{hence}\quad S(\min(f,g))\le\min(Sf,Sg).
\end{equation}
The minimum is componentwise, so the minimizing branch can vary
with phase and class.  Applying this to the two exponentials
proves that $C^j/V$ is a free supersolution on the entire line,
in all four phases.

At a stored coordinate where a row is required, each actual
successor value is at most $\widehat E/V$: equality holds inside,
and outside $C^j/V\le\lceil C^j\rceil/V$.
Multiplying this by the nonnegative coefficients and their
upper ceilings, \eqref{hl:linerow} divided by $DV$ proves
\eqref{hl:lineexact}.  This includes the rows at $70$ and $-70$;
a step parallel to the other axis has coordinate displacement zero.

At $|k|\ge71$, every neighbour coordinate is either outside the
stored interval or is $70$ or $-70$.  Outside, $\ell^j=C^j/V$;
at the possible stored neighbour the floor cap in
\eqref{hl:linecaps} gives $\ell^j\le C^j/V$.  Thus
\[
 \sum_e t_e\ell^j_{p_e,q_e}(k+(d_e)_j)
 \le\sum_e t_e C^j_{p_e,q_e}(k+(d_e)_j)/V
 \le C^j_{p,q}(k)/V=\ell^j_{p,q}(k).
\]
This proves the first outside rows $k=\pm71$ as well as all
farther ones; it explains why an inequality on only the inside
outgoing row would not suffice.

Finally the stored C,D bounds are \eqref{hl:linecaps}.  If
$k\ge71$, choose the decaying positive branch, obtaining
$C^j_{p,q}(k)/V\le Z_{\max}r^{71}/V\le1$ by
\eqref{hl:tailbound}; if $k\le-71$ use the negative branch.
This actually bounds all phases in these tails by one.  All
values are finite and nonnegative because the data and every
fixed rational power of $r$ are finite and nonnegative.
\end{proof}

\begin{lemma}[Lines in the plane]\label{hl:planelemma}
The function $F^v_\infty$ satisfies
\[
 \mathcal WF^v_\infty(p,q,x)\le F^v_\infty(p,q,x)
        \quad(x\ne v),\qquad
 F^v_\infty(p,q,v)\ge g_p.
\]
Consequently $g+T_{X,d_0}F^v_\infty\le F^v_\infty$ for each case.
The free inequality is not asserted at the target.
\end{lemma}

\begin{proof}
Lift the $j$th line to $f_j(p,q,x)=\ell^j_{p,q}(x_j-v_j)$.
For each free edge its next coordinate is exactly
$x_j-v_j+(d_e)_j$.  Therefore its two-dimensional free row
\eqref{hl:free} is exactly the left side of
\eqref{hl:lineexact}.  In particular, edges perpendicular to
that axis retain their multiplicity, displacement zero in that
coordinate, and their next phase and class; they are not omitted.
This proves the free row for $f_j$ in A,B everywhere and in C,D
whenever $x_j-v_j\ne0$.

At an A,B state choose any line attaining the minimum in
\eqref{hl:Finfty}.  By \eqref{hl:min},
$\mathcal WF^v_\infty\le\mathcal Wf_j\le f_j=F^v_\infty$ there.
At a C,D state with $x\ne v$, if a minimizing line has nonzero
coordinate the same argument applies.  Otherwise its value is
one.  The other coordinate is nonzero and its line value is at
most one by Lemma~\ref{hl:taillemma}.  Since the zero-coordinate
line was minimizing, this other value must equal one as well.
It is therefore also a minimizing line and its row is available.
This proves the asserted free inequality at every non-target state.

At $x=v$ the two C,D values are exactly one, and the two A,B
values are nonnegative.  Thus $F^v_\infty\ge g$ there.
At $c$ and $v$ the stationary row is zero.  At every other site
$g=0$, so \eqref{hl:entrywise} and the proved free inequality give
$TF^v_\infty\le F^v_\infty$.  At $c$ the reward is zero; at $v$
the terminal domination just proved applies.  This establishes
the stopped supersolution on the whole state space.
\end{proof}

\begin{lemma}[Patching without a boundary gap]\label{hl:patchlemma}
For each case the function $F$ in \eqref{hl:Fpatch} satisfies
\[
 0\le F\le F^v_\infty,\qquad F=g\text{ at }c,v,
                  \qquad g+TF\le F
\]
on the entire lattice.
\end{lemma}

\begin{proof}
Nonnegativity, the cap, and the terminal values follow from
\eqref{hl:patchcaps}.  At an inside nonterminal state the exact
successor value is $\widehat U/V$, including every outside
successor.  Divide \eqref{hl:patchrow} by $DV$ and use the
coefficient ceilings to obtain $TF\le F$.  This proves every
inside boundary row without identifying opposite sides or
discarding an outgoing edge.  At an outside state, monotonicity
and Lemma~\ref{hl:planelemma} give
\[
 TF\le TF^v_\infty\le F^v_\infty=F.
\]
This includes rows entering the patch, since the cap holds at
every possible inside successor.  Both terminals are inside $P$,
so outside $g=0$.  At either terminal $TF=0$ and $F=g$, proving
the assertion everywhere.
\end{proof}

\subsection{True suffixes and exact prefixes}

\begin{lemma}[Finite horizons]\label{hl:bellman}
For each case, phase, class, position, and integer $n\ge0$,
\[
 \sum_{j=0}^{n}T^jg\le F.
\]
The left side counts the relaxed first-hit paths of length at
most $n$ and terminal reward $g$.  Their all-length sum is also
bounded by $F$.
\end{lemma}

\begin{proof}
There are finitely many outgoing directions and phases at each
state, so each finite-horizon value is a finite sum.  For $n=0$,
$g\le g+TF\le F$.  If $H_n=\sum_{j=0}^nT^jg\le F$, then
$H_{n+1}=g+TH_n\le g+TF\le F$ by positivity and
Lemma~\ref{hl:patchlemma}.  This is the induction.  Expanding
$T^jg$ counts every path once with its edge product and terminal
reward.  A path hitting a terminal earlier has a zero outgoing
factor.  Thus nonzero terms are precisely the first-hit paths
with nonzero terminal reward.  Taking the increasing supremum
over $n$ proves the all-length statement.  In particular no
stationary uniqueness or convergence rate is needed.
\end{proof}

We next check that this relaxed bound applies to the paper's
true suffixes with their full history.  This is also a direct
proof of the part of Appendix C's stationary-domination lemma
used here; changing the potential does not change that lemma.

\begin{lemma}[Permanent initial data and true suffix domination]
\label{hl:suffix}
Fix a true length-twelve prefix in $\mathcal P_{X,d_0}$.
For a specified phase at its endpoint, the total weight of its
true continuations to their first hit of $v$, with the actual
past retained, is at most $F(p,q(P),x(P))$.  This holds first
at every finite suffix horizon and then for their supremum.
\end{lemma}

\begin{proof}
At a suffix departure time $j\ge12$, the previous ten steps
form a raw window in $\Omega_{10}$.  The next true bond is
not used in that window and cannot reverse the last step, so
the direction is admitted.  Its new class is precisely the
quotient successor.  The true trail avoids $c$ and stops on its
first hit of $v$, so it satisfies the two terminal rules.

Write $a=u+d_0$.  The true initial bond $\{u,a\}$ was used at
time zero and can never be used again in either direction.
Both known sites have been departed from, at times zero and
one respectively.  These departures exist in a surviving
length-twelve prefix and neither site is $v$.  At every suffix
time the remembered departure indices are $j-10,\ldots,j-1$,
all at least two.  Thus the known initial departures are
\emph{outside} the window for the whole suffix, not merely at
the split.  If the memory detects a revisit at either known
site, it detects a further departure in addition to its initial
departure and the proposed current departure.  This would be a
third departure, which is impossible: at $u$ it would need at
least five distinct incident bonds, whereas only three avoid
$c$; at $a$ its initial arrival/departure and two additional
arrival/departure pairs would need six among four bonds.
Here $a$ is one of $(-1,1),(-2,0),(-1,-1)$, distinct from
$u,c,v$ in each case.  The killing rule therefore kills no
true continuation.  Every remaining true departure at a known
site is a revisit, with exactly the retained $1/5$ diagonal
weight and zero response.

At any other site, a memory-detected revisit is a true revisit,
so the weights agree.  A true revisit forgotten by memory has
weight $1/5$ in the same phase and no response; choose that
diagonal edge in the relaxed history, whose coefficient is at
least $1/3$.  At a true first departure, memory also sees a
first departure, with the same incoming/outgoing directions
and the same diagonal coefficient.  If the phase responds,
only $c$ is kept dead in the relaxed reward.  When the other
dead neighbours have been forgotten, an existing sole pair's
coefficient can increase from $b'_\tau$ to $b_\tau$, and the
remaining added pairs have nonnegative coefficients.  If no
true pair survives the true reward is zero.  Since the
outgoing true neighbour is live and the two $h$ entries are
positive, these cases prove that each true responding weight
is at most its relaxed weight.

Map a true suffix and its phase history to this relaxed
position/phase path.  Its successive positions determine
every absolute direction, and its phases determine the phase
history.  The map is injective for a fixed prefix and initial
phase even though window classes have been identified.  The
edge comparisons above therefore bound the sum over true
histories at any finite horizon by the relaxed stopped sum.
Lemma~\ref{hl:bellman} bounds the latter by $F$.  Taking
suprema proves the all-length assertion.  The finite-horizon
argument retains all true prefix history; it does not assert
that the quotient decides actual dead sites or actual revisits.
\end{proof}

\begin{lemma}[The four phases realize the long ordering]
\label{hl:phaseidentity}
For a true trail $w$ of length $n>12$, its contribution to
$(M_1h)_X$ is exactly the sum of its phase-history products
starting with one unit in each of A and B and terminal reward
$(0,0,1,1)$, using its actual $\kappa_i$ in the phase changes.
Splitting any such history after twelve departures factors its
product into the exact prefix factor and true suffix factor.
\end{lemma}

\begin{proof}
A nonzero rewarded history starting in A changes exactly once,
from A to C, at a turning first departure $i$.  Its product is
\[
 3^{-n}(3/5)^{\operatorname{rev}(w)}\kappa_i
 x_{\rm s}^{\#\{\text{all straight first departures}\}}
 x_{\rm t}^{\#\{\text{turning first departures at indices }j\le i\}}.
\]
This contains exactly the $x$ factors up to that turning
departure in the straight-first ordering.  A rewarded history
starting in B changes exactly once, B to D, at a straight
first departure $i$.  Its product is
\[
 3^{-n}(3/5)^{\operatorname{rev}(w)}\kappa_i
 x_{\rm s}^{\#\{\text{straight first departures at indices }j\le i\}},
\]
again exactly the corresponding ordered kernel term.  There
are no other rewarded histories.  Zero rewards give zero
terms.  Summing gives \eqref{hl:kernel} contracted with $h$.
Each edge product factors at twelve, and grouping the prefix
histories by their phase gives exactly \eqref{hl:phaseupdate}.
This factorization is algebraic; no independence of suffix
weights from the full true past is assumed.
\end{proof}

\begin{theorem}[Finite-to-global implication]\label{hl:theorem}
H3-line implies, for the true-trail kernel
\eqref{hl:kernel} with exactly \eqref{hl:order}, that
\begin{equation}\label{hl:result}
  M_1h\le\rho_0h,\qquad h=(1,6/5),\quad\rho_0=967725/10^6.
\end{equation}
The same bound holds for every finite true-trail length cutoff,
as required by (H3$_h$) in Section 7.
\end{theorem}

\begin{proof}
Fix a finite cutoff $R_0$ and a first-direction case.  Trails of
length at most $\min\{R_0,12\}$ contribute a subsum of
$S_{X,d_0}$, because their ordering is exactly the short
ordering used there.  A longer trail has a unique surviving
length-twelve prefix.  By Lemma~\ref{hl:phaseidentity}, split
its phase histories at twelve and group by the phase at that
point.  Lemma~\ref{hl:suffix} bounds the true suffix sum of
length at most $R_0-12$ by the corresponding $F$ value.
Multiplying by its nonnegative exact prefix weight and summing
over the finitely many prefixes gives \eqref{hl:J} as an upper
bound.  If $R_0\le12$, the nonnegative prefix term can simply
be added to the short bound.  This covers $R_0=0$ also.
Coincident quotient/endpoint states do not remove any
multiplicity: every distinct prefix still appears separately
with its own weight.

For row B the reflection $(x_1,x_2)\mapsto(x_1,-x_2)$ fixes
$c,u,v$ and the incoming W direction.  It interchanges initial
N and S and is a bijection on true trails at every length.
It preserves all used bonds, departure counts, dead flags,
straight/turn/revisit kinds, opposite/perpendicular pair types,
and graph distances defining the shifts.  It also preserves
the path indices.  Thus it preserves both the short rational
keys, including their index tie rule, and the long
straight-first order.  Every kernel term agrees.  The B/N
and B/S true sums are therefore equal.  This needs no symmetry
of the quotient numbering or the arrays.  For C keep all
three first-direction cases.  Initial E enters $c$ and is
forbidden.  We obtain, at every cutoff,
\begin{align*}
 (M_1h)_B&\le2J_{B,\mathrm N}+J_{B,\mathrm W}
 \le\frac{967724846}{10^9}
 <\frac{967725000}{10^9}=\rho_0h_B,\\
 (M_1h)_C&\le J_{C,\mathrm N}+J_{C,\mathrm W}+J_{C,\mathrm S}
 \le\frac{1159917640}{10^9}
 <\frac{1161270000}{10^9}=\rho_0h_C.
\end{align*}
The integer differences in the strict comparisons are $154$
and $1352360$, respectively.  Here the notation on the left
initially denotes the cutoff sum.  Taking its increasing
supremum gives the full nonnegative series.  Regrouping by
$Y,s$ then gives precisely $M_1h$ in \eqref{hl:kernel}.
This proves its finiteness and \eqref{hl:result}, without a
prior summability assumption.
\end{proof}

\subsection{Literal data contract and exact replay}

The source directory for the literals below is
\texttt{/home/ubuntu/aklt/reports/MATH\_FIT/data}.
The following description fixes their mathematical meaning;
generation histories and fixed-point iteration logs are not
part of the hypothesis.  All binary integers are little-endian.
All array payload values are unsigned 32-bit integers.  Header
integers are signed 64-bit integers.  Exact lengths, headers,
and absence of trailing payload are checked.

The text file \texttt{windows\_m10\_min.txt} starts with
$(60092,3741)$, followed by $60092$ pairs $(e(\omega),\pi(\omega))$
and then $3741$ rows consisting of a representative code,
four successors and four kinds, in E,N,W,S order.
It is verified by the raw reconstruction preceding
\eqref{hl:windowcounts}.

The file \texttt{axis\_z\_integer.bin} contains $16\cdot3741$
values, no header, in order $\sigma,p,q$, with
$\sigma=+e_1,-e_1,+e_2,-e_2$ and $p=\mathrm A,\mathrm B,\mathrm C,\mathrm D$.
The certified extrema and a valid uniform residual are
\begin{equation}\label{hl:eta}
 Z_{\min}=10486,\qquad Z_{\max}=83704870,\qquad
 \eta=2951081685866839636.
\end{equation}
Every row residual is at least this $\eta$; the replay computes
the minimum over all rows.  The tail comparison is exactly
$83704870\cdot12^{71}\le2^{20}\cdot13^{71}$, whose ratio
of the two sides is approximately $0.271670<1$.

The file \texttt{line\_integer.bin} has header
$[10,70,3741,20,48]$ and then the $E$ array in order
$j,p,q,k$, with $j=1,2$ and $k=-70,\ldots,70$.
Its zero-based index is
\[
 (((j-1)4+p)3741+q)141+(k+70),\quad p=0,1,2,3.
\]
The files \texttt{line\_init\_integer.bin} and
\texttt{line\_tails\_integer.bin} have no header.  The former
contains the derived cap floors $A$ in the same order; the
latter contains the derived cap ceilings $B$ in order
$j,p,q,-,+$.  The replay recomputes these numbers from $Z$
and exact powers of $12/13$, comparing every entry.  These
auxiliary files can be omitted by a checker that regenerates
the tables directly; their values are not independent witnesses.

The five patch files are
\[
 \texttt{line\_}X\texttt{\_m10\_R20\_d}j\texttt{\_integer.bin.gz},
 \qquad (X,j)=(B,1),(B,2),(C,1),(C,2),(C,3).
\]
Their decompressed header is
$[10,12,20,3741,1681,v_1,v_2,20,j]$.
It is followed first by $4\cdot3741$ values copying
$E^1_{p,q,0}$, which the replay checks entry by entry, and
then by $4\cdot3741\cdot1681$ patch values $U$ in order
$p,q,x$, with spatial index $(x_1+20)41+(x_2+20)$.
The copied line-zero vector is redundant metadata, \emph{not}
the old constant exterior cap of Appendix C; it must not be
used to extend the patch.  There are no further values.

\begin{center}\small
\begin{tabular}{lrrrr}
\toprule
Object&Values&Required row inequalities&Maximum value&Bits\\
\midrule
Four tilted vectors&59856&59856&83704870&27\\
Two stored lines&4219848&4204884&10768774&24\\
One patch&25154484&25124556&1066965\text{ (at most)}&21\\
Five patches&125772420&125622780&1066965&21\\
\bottomrule
\end{tabular}
\end{center}
Every line value is also subjected to its cap check; all
$14964$ C,D zero-coordinate values and $2094960$ C,D
nonzero-coordinate values are checked as specified.
The line row count is $2\cdot4\cdot3741\cdot141-2\cdot2\cdot3741$.
Each patch has $29928$ terminal values and $2394240$ rows
on its spatial boundary; these are included in the displayed
counts.  The main row count is
$59856+4204884+125622780=129887520$, before cap,
terminal, indexing, and prefix checks.

The exact prefix enumeration gives, in the same five-case order,
\[
\begin{array}{c|rrrrr}
 &B/\mathrm N&B/\mathrm W&C/\mathrm N&C/\mathrm W&C/\mathrm S\\\hline
 \text{nodes}&139985&161929&105025&155673&145485\\
 \text{surviving prefixes}&89150&103257&66947&99168&92748\\
 \text{short first hits}&937&922&715&1036&861\\
 \text{bits of }A_{X,d_0}&422&422&421&422&422
\end{array}
\]
Counts are diagnostics in addition to the rational comparisons
\eqref{hl:clearedprefix}; counts alone do not prove those
comparisons.  The exact rational values are printed in
\texttt{replay.log}.  Their B row and C row divided by $h_C$
are approximately $0.9677248435916523$ and
$0.9665980319004569$, respectively.  The proof uses the
integer ceilings in \eqref{hl:prefixceil} instead of decimals.

The supplied \texttt{replay.cpp} is an implementation of this
finite statement: it reconstructs all raw words and verifies
the quotient, regenerates rational coefficients and both kinds
of ceilings, verifies every line/patch/interface inequality,
and enumerates the prefixes with ungrouped exact short sorting.
It reads these literals but includes or calls no checker from
\texttt{MATH\_FIT/programs}.  Coefficient and row arithmetic
uses exact integers, and cap and prefix arithmetic uses
arbitrary-precision rationals/integers.  Fixed-width row sums
use unsigned 128-bit integers: every input value is at most
$2^{32}-1$.  The two free rewards are $9918938/10^7$ and
$9666078/10^7$, both less than one, and every reduced reward
is no larger.  Thus every coefficient, even after tilting,
is at most $(x_{\rm t}/3)(13/12)=650/1683<1/2$;
its rounded numerator is at most $D/2<2^{48}$.
A row has at most six individual edges, so
its sum is less than $2^{83}$.  There is no possible overflow
in these rows.  Prefix accumulators use arbitrary precision.
The maximum observed sum widths are $75$ bits (tilts), $72$
bits (lines), and $69$ bits (patches).

As additional checks of the delicate interfaces, the replay
directly evaluates the exact, unrounded free rows at
$k=-100,-71,71,100$ on both axes for every phase and class:
all $119712$ inequalities hold, including the rows at $\pm71$
whose successors can be stored values at $\pm70$.
It also evaluates all $14964$ C,D rows at coordinate zero,
which the certificate does not require.  For these particular
data even those rows hold; their largest exact ratio of left
side to value is $328000259/377880576<1$, at axis 1,
phase C, class 0.  This stronger observed property is not
used in the minimizing-line argument.  The finite tail probes
are supplementary tests, not a replacement for the proof at
all integer distances.

Running \texttt{bash replay.sh} builds and executes that replay.
The accompanying log records successful checks of the displayed
finite assertions.  This is an exact arithmetic replay and a
written soundness proof; it is not a claim of a new Lean-kernel
verification.  The hypothesis and every finite assertion remain
the explicit mathematical ones above, not a hash or exit status.

\subsection{Optional deduction of (H2) from the same vectors}

The following three additional integer comparisons are satisfied
by \eqref{hl:eta}:
\begin{equation}\label{hl:foldedtests}
 8000\eta\ge D Z_{\max},\qquad
 Z_{\max}<8000 Z_{\min},\qquad
 1600000^{10}<2\cdot1495813^{10}.
\end{equation}
They are checked by the replay.  They are unnecessary if the
original memory-six (H2) input is retained.

\begin{proposition}\label{hl:H2}
H3-line together with \eqref{hl:foldedtests} implies the paper's
open-trail input (H2) with $q=7999/8000$ and $C_V=16000$.
\end{proposition}

\begin{proof}
Let $\gamma(i,d)$ be the A-phase diagonal weight at class $i$, namely
$x_{\rm s}/3,x_{\rm t}/3,1/5$ according to kind, and define
$A_\pm f(i)=\sum_d\gamma(i,d)r^{\pm(d)_1}f(id)$.
Discarding the nonnegative A$\to$C edges from
\eqref{hl:tiltcheck}, dividing by $D$, and using the exact
coefficient bounds gives
\[
 A_\pm Z^{\pm e_1}_{\mathrm A,\cdot}
 \le Z^{\pm e_1}_{\mathrm A,\cdot}-\eta/D
 \le (7999/8000)Z^{\pm e_1}_{\mathrm A,\cdot}.
\]
The last step uses $\eta/D\ge Z_{\max}/8000$.
Put $v_i=\sqrt{Z^{+e_1}_{\mathrm A,i}Z^{-e_1}_{\mathrm A,i}}>0$
for a class $i\in Q$.  For the untilted A diagonal matrix
$A_0$, finite-sum Cauchy--Schwarz yields
\[
\begin{split}
 (A_0v)_i
 &=\sum_d\sqrt{\gamma(i,d)r^{(d)_1}Z^{+e_1}_{\mathrm A,id}}
             \sqrt{\gamma(i,d)r^{-(d)_1}Z^{-e_1}_{\mathrm A,id}}\\
 &\le\sqrt{(A_+Z^{+e_1}_{\mathrm A})_i
                 (A_-Z^{-e_1}_{\mathrm A})_i}
 \le qv_i,\qquad q=7999/8000.
\end{split}
\]
Also $\max v/\min v\le Z_{\max}/Z_{\min}<8000$.
Let $\widehat A_0$ be the raw diagonal memory-ten matrix.
The verified quotient gives
$\widehat A_0(\pi^*v)=\pi^*(A_0v)\le q\pi^*v$;
surjectivity preserves the extrema of $v$.

For clarity, the open-trail argument works for this possibly
irrational vector.  Let $a_n(\omega)$ be the true total
open-trail weight with last window $\omega$, for $n\ge10$,
and $W_n$ its total weight, with any fixed start/incoming
direction and initial reversal forbidden.  A true extension
is admitted by its window.  A detected revisit has exactly
weight $1/5$; a forgotten true revisit has weight $1/5$ no
greater than either relaxed first-departure weight.  First
departures otherwise agree.  The injective map to prefix and
last direction therefore gives $a_{n+1}\le a_n\widehat A_0$.
Iteration with the positive raw vector $\pi^*v$ proves
\[
 W_{10+j}\le(\max v/\min v)q^jW_{10}.
\]
At each true prefix there are at most three next directions,
including at the initial step because reversal is forbidden,
and each has weight at most $x_{\rm t}/3$.  Consequently
$W_t\le x_{\rm t}^t$ for every $t\ge0$.  Since
$x_{\rm t}/q>1$, the constants for lengths below ten and
for lengths at least ten are bounded respectively by
$(x_{\rm t}/q)^{10}$ and
$8000(x_{\rm t}/q)^{10}<16000$.
The last inequality is the final comparison in
\eqref{hl:foldedtests}, since $187\cdot7999=1495813$.
Thus $W_n\le16000q^n$ at all lengths, uniformly in the
start and incoming direction.  This is exactly (H2).
\end{proof}

\subsection{The modified Section 7 closing inequalities}

Assume (H1) at $961/1000,187/200$, and either retain the
original (H2) constants or use Proposition~\ref{hl:H2}.
The exposure/response recursion in Section 7 allows arbitrary
fixed orderings and uses the same type-dependent coefficients.
It therefore applies to \eqref{hl:order}, \eqref{hl:kernel}.
Here we give all constants in its witness shift and renewal.
Use
\begin{equation}\label{hl:closingconstants}
 \bar\alpha=\frac{18094}{100000},\qquad
 \bar\beta=\frac{36751}{100000},\qquad
 C=\bar\alpha+2\bar\beta=\frac{22899}{25000},\qquad H=\frac65.
\end{equation}
For either chosen $q,C_V$, put
\begin{equation}\label{hl:shiftconstants}
\begin{gathered}
 b=\frac{1+q}{2},\quad r_0=\frac qb,\quad
 B_*=2HCC_V b^{-2}\frac{r_0}{(1-r_0)^3}
             =\frac{16HCC_Vq}{(1-q)^3},\\
 \varepsilon=\min\left\{\frac{1-q}{4},\frac{1291}{80000B_*}\right\},
 \qquad\theta=1-\varepsilon,\qquad
 \rho=\frac{1+\rho_0}{2}=\frac{78709}{80000},\\
 C_T=2(12/5)^{24}\left(\frac{C_V}{1-q}\right)^2,\qquad
 A=\frac65\frac{\max\{2,C_T\}}{1-\rho}
   =\frac{96000}{1291}C_T.
\end{gathered}
\end{equation}
The last equality uses $C_T>2$, valid for both choices below.
The exact values are
\begin{equation}\label{hl:closingtable}
\begin{array}{c|c|c}
 &\text{retained (H2)}&\text{folded (H2)}\\\hline
 q&1119948287353/1120437614225&7999/8000\\
 C_V&10&16000\\
 \varepsilon&
 \displaystyle\frac{295429241889217285418799689}
 {38634126128180928127866293784580458800250}&
 \displaystyle\frac{1291}{11524003395010560000000}\\[6pt]
 C_T&\displaystyle2(12/5)^{24}
                  (5602188071125/244663436)^2&
       \displaystyle2(12/5)^{24}(128000000)^2\\
 A&(96000/1291)C_T&(96000/1291)C_T
\end{array}
\end{equation}
In both columns the second candidate in the minimum for
$\varepsilon$ is the smaller; these comparisons, the fractions,
and $C_T>2$ are checked by exact rational arithmetic in the
replay.  Approximately, the retained column has
$\varepsilon=7.64685\cdot10^{-15}$,
$C_T=1.39855\cdot10^{18}$, $A=1.03997\cdot10^{20}$;
the folded column has $\varepsilon=1.12027\cdot10^{-19}$,
$C_T=4.37039\cdot10^{25}$, $A=3.24986\cdot10^{27}$.
These approximations are not used as inequalities.

Here are the summability and shift bounds explicitly.
A response pair of a length-$n$ trail has shift at most
$n/2+1$: at departure $i$ its site is within $i$ of $u$ and
within $n-i$ of $v$, and each response vertex is its neighbour.
There are at most $n$ first departures and total response
coefficient at most $C$ at each.  All $x_i\ge1$, so every
ordered partial product is at most the full product.  With
(H2), each entry of $M_\theta$ is therefore bounded by
\begin{equation}\label{hl:entries}
 CC_V\sum_{n\ge1}nq^n\theta^{-n/2-1}
 =CC_V\theta^{-1}\frac{q/\sqrt\theta}{(1-q/\sqrt\theta)^2}<\infty
       \quad(q^2<\theta\le1).
\end{equation}
Our choice gives $q^2<q<b<\theta<1$.  For any integer
$0\le s\le n/2+1\le n+1$, the mean value theorem and
$\theta>b$ give
\[
 0\le\theta^{-s}-1\le\varepsilon s b^{-s-1}
                       \le\varepsilon(n+1)b^{-n-2}.
\]
On multiplying each pair reward by $h_Y/h_X\le H$ and
using (H2), this yields
\begin{equation}\label{hl:shiftineq}
 \frac{((M_\theta-M_1)h)_X}{h_X}
 \le\varepsilon HCC_V b^{-2}\sum_{n\ge1}n(n+1)r_0^n
 =\varepsilon B_*\le\frac{1291}{80000}.
\end{equation}
The identity in the sum is
$\sum_{n\ge1}n(n+1)r_0^n=2r_0/(1-r_0)^3$.
Together with Theorem~\ref{hl:theorem},
\begin{equation}\label{hl:shifted}
 M_\theta h\le(38709/40000+1291/80000)h
                   =(78709/80000)h=\rho h.
\end{equation}
Thus the Section 7 pair envelopes, denoted here by
$\mathcal F_X(d)$ to distinguish them from the Bellman
potential, obey exactly
\begin{equation}\label{hl:recursion}
 \mathcal F_X(d)\le\sum_{Y\in\{B,C\}}\sum_{s\ge0}
     M_{XY,s}\mathcal F_Y(d-s)+C_Tq^{2d}\qquad(d\ge1),
\end{equation}
with $0\le\mathcal F_X\le2$ and constant extension for $d<0$.
The factor $(12/5)^{24}$ and the count of twelve exceptional
U departures in the second trace are unchanged.

To see the renewal bound directly, let
$(Kf)_X(d)=\sum_{Y,s}M_{XY,s}f_Y(d-s)$ and
$G_X(d)=2$ for $d\le0$, $C_Tq^{2d}$ for $d\ge1$.
Then $\mathcal F\le K\mathcal F+G$ on all integers, and
\eqref{hl:shifted} implies $K^kh\le\rho^kh$ when $h$ is
constant in $d$.  Also, for $e_X(d)=h_X\theta^d$,
$Ke\le\rho e$.  Since $\min h=1$,
$\mathcal F\le2h$ and $G\le\max\{2,C_T\}e$.
Finite iteration and positivity give
\[
 \mathcal F_X(d)\le2\rho^kh_X+
 \max\{2,C_T\}h_X\theta^d\frac{1-\rho^k}{1-\rho}.
\]
The unshifted and shifted row bounds justify every application
and every nonnegative rearrangement, including negative $d$.
Letting $k\to\infty$ proves
\begin{equation}\label{hl:decay}
 \mathcal F_X(d)\le A\theta^d\quad(d\ge0),\qquad A\ge2,
\end{equation}
with the explicit constants above.  This is the pair property
$\mathsf P(A,\theta)$ used in Section 5, by the same
ordinary-domain specialization in Section 7.

For reference the later choices can also remain literally
those of the paper: take $l_0=64$ and
\[
 E_0=1+4A+10A\theta^{-4}
       \left[48\theta^{14}\left(17+\frac1{1-\theta}\right)+1\right],
 \quad p_0=e^{-2E_0}>0,
\]
\[
 K_\theta=36\theta^{-1}10!\left(\frac{12}{-\log\theta}\right)^{10},
 \qquad M_\star=\left\lceil\max\{390^4,K_\theta^2,1\}\right\rceil.
\]
These are finite.  The Section 4 constant is then
$\delta=p_0\alpha^{1024}>0$, with the paper's unchanged
frame finite-energy $\alpha$.  The ensuing choice of one
sufficiently large square and the spectral-gap argument need
no further finite trail certificate.

\subsection{Scope of the replacement}

The following specifies precisely how this section relates to
the present Appendix C, using its source labels to avoid any
dependence on theorem numbering.
\begin{itemize}
\item The raw memory-automaton definition, quotient lemma
\texttt{l:automaton-quotient}, and quotient reconstruction
\texttt{p:automaton10} are unchanged.  All original H2
definitions, tests, symmetry reductions, and the open-trail
lemma \texttt{l:alm} remain valid if H2 is retained.  If it is
folded, Proposition~\ref{hl:H2} replaces the need for the
memory-six numerical input and its $C_V=10$ corollary;
the general lemmas are still valid.
\item The phase-weight definition \texttt{d:phaseweights},
phase-expansion lemma \texttt{l:phase-expansion}, stationary
walk \texttt{d:stationary-walk}, and stationary domination
\texttt{l:stationary-domination} remain valid with the newly
specified $h_C=6/5$ in every reward.  Their proofs use only
positivity, $b'_\tau\le b_\tau$, and the indicated long
ordering, and are proved again here where needed.  The
Bellman lemma \texttt{l:Bellman-soundness} is unchanged.
\item Replace \texttt{d:stationary-certificate}, its
dyadic-$h$ response-ceiling specialization, and
\texttt{l:stationary-local-global} by Definition~\ref{hl:finite}
and Lemmas~\ref{hl:taillemma}--\ref{hl:patchlemma}.
The old lemma itself remains true for its old constant-cap
certificates, but those are no longer the input used here.
In particular the old $V=2^{16}$, radius-70 spatial arrays
and constant exterior $Z/V$ cannot be reused as the new
array contract.
\item The mathematical definition \texttt{d:prefix-bound},
its exhaustive recursion and short sorting, and
\texttt{l:prefix-domination} remain valid with the new $h$,
the new $F$ and $V$, and the same cutoff $12$.
Replace the old numerical denominator by
\eqref{hl:Delta}, the full spatial lookup by the radius-20
lookup specified here, and every numerical prefix
evaluation by \eqref{hl:prefixceil}.  The enumeration counts
are unchanged because they depend only on trails, not on
$h$ or the potential.
\item The kernel reflection lemma
\texttt{l:kernel-\allowbreak reflection} is unchanged.  Replace
\texttt{p:stationary-\allowbreak check} and its binary/indexing contract
by the literal and finite contracts here, and replace the
numerical part of \texttt{l:stationary-\allowbreak assembly} by
Theorem~\ref{hl:theorem}.  Its split-by-first-direction
argument remains valid.
\end{itemize}
In Section 8 the H3 assertion and its last proof paragraph
must cite H3-line with $V=2^{20}$, the two one-dimensional
radius-70 arrays and the five radius-20 patches, and use
$h=(1,6/5)$ and $\rho_0=967725/10^6$.  In Section 7 change
the kernel-order constants and the numerical witness shift
as above; its response derivation, recursion and renewal
proofs are otherwise valid.  H1 and Appendix D can be kept
as written.  None of the finite checks in this section
establishes the separate H1-tree proposal in the report.

The sketch in the report needs the following qualifications
to be a complete certificate.  Floors are upper-cap tests
on stored line values and ceilings are used on outside
successors; the infinite tail itself is the exact rational
minimum of exponentials.  C,D zero-coordinate lines need
no free row, and the free inequality at the target is neither
asserted nor needed.  The terminal reward is dealt with by
zero terminal rows.  Both patch boundary directions must be
handled with the global cap and with all inside boundary
rows.  The known initial sites and edge are permanent and
depend on the first direction; they cannot be dropped or
replaced by an unsplit potential.  All $h$-dependent
coefficients, ordering keys, and denominators must be
recomputed for exact $6/5$.  The proof above and the replay
implement these qualifications without changing any witness
array or any claimed contraction constant.
